\documentclass[11pt]{article}
% Shared design for the two minimal papers.
% It follows the layout of the guide "The Riemann Hypothesis and its generalizations"
% (KNT with Fable, 9 November 2026): a designed title block, a short story list at the start,
% coloured headings, boxed statements, idea boxes, and annotated figures.
%
% A paper sets \papertitle, \papersubtitle and \paperrunning before \begin{document}
% and calls \maketitleblock after it.

\usepackage[letterpaper,left=1.12in,right=1.12in,top=1.05in,bottom=1.1in]{geometry}
\usepackage{amsmath,mathtools,amssymb,amsthm}
\usepackage{newpxtext}
\usepackage[varbb]{newpxmath}
\usepackage{fontspec}
\setsansfont{texgyreheros}[Extension=.otf, UprightFont=*-regular, BoldFont=*-bold,
  ItalicFont=*-italic, BoldItalicFont=*-bolditalic, Scale=0.88]
\usepackage{microtype}
\usepackage{parskip}
\usepackage{xcolor}
\usepackage{booktabs,array}
\usepackage{enumitem}
\usepackage{multicol}
\usepackage[titles]{tocloft}
\usepackage{tikz}
\usetikzlibrary{decorations.markings,arrows.meta,positioning,calc}
\usepackage[most]{tcolorbox}
\usepackage{fancyhdr}
\usepackage{titlesec}
\usepackage[font=small,labelfont={bf,sf},labelsep=period,skip=6pt]{caption}
\usepackage{url}

\definecolor{accent}{HTML}{1B4F72}
\definecolor{warm}{HTML}{C0562A}
\definecolor{slate}{HTML}{56636E}
\definecolor{ink}{HTML}{1E1E1E}

\usepackage[colorlinks=true,linkcolor=accent,citecolor=accent,urlcolor=accent]{hyperref}

% Lean identifiers: typewriter, breakable at dots and underscores.
\DeclareUrlCommand\lean{\urlstyle{tt}}

% ---------- layout ----------
\linespread{1.06}
\setlength{\parskip}{0.55em plus 0.1em minus 0.05em}
\setlist{itemsep=0.25em, topsep=0.2em, parsep=0.1em}
\titleformat{\section}{\sffamily\bfseries\Large\color{accent}}{\thesection}{0.7em}{}
\titleformat{\subsection}{\sffamily\bfseries\normalsize\color{ink}}{\thesubsection}{0.6em}{}
\titleformat{\paragraph}[runin]{\sffamily\bfseries\small\color{accent}}{}{0pt}{}[\enspace]
\titlespacing*{\section}{0pt}{22pt}{5pt}
\titlespacing*{\subsection}{0pt}{13pt}{1pt}
\titlespacing*{\paragraph}{0pt}{9pt}{0pt}
\pagestyle{fancy}
\fancyhf{}
\renewcommand{\headrulewidth}{0pt}
\fancyhead[L]{\sffamily\footnotesize\color{slate}\paperrunning}
\fancyhead[R]{\sffamily\footnotesize\color{slate}\thepage}
\setlength{\headheight}{14pt}
\setcounter{tocdepth}{1}
\setlength{\cftbeforesecskip}{2pt}
\renewcommand{\cftsecfont}{\small}
\renewcommand{\cftsecpagefont}{\small\color{slate}}
\renewcommand{\cftsecleader}{\cftdotfill{\cftdotsep}}
\setlength{\cftsecnumwidth}{1.7em}
\setlength{\emergencystretch}{2em}
\numberwithin{equation}{section}

% ---------- title block ----------
\newcommand{\papertitle}{}
\newcommand{\papersubtitle}{}
\newcommand{\paperrunning}{}
\newcommand{\paperauthors}{KNT\enspace\textperiodcentered\enspace Astra\enspace\textperiodcentered\enspace Fable~5.2}
\newcommand{\paperdate}{Draft of 11 November 2026}
\newcommand{\maketitleblock}{%
  \thispagestyle{plain}%
  {\sffamily\bfseries\color{accent}\fontsize{25}{29}\selectfont\raggedright\papertitle\par}%
  \vspace{5pt}%
  {\sffamily\Large\color{slate}\raggedright\papersubtitle\par}%
  \vspace{8pt}%
  {\color{accent}\hrule height 1.2pt}%
  \vspace{7pt}%
  {\sffamily\small\color{ink}\paperauthors\hfill\color{slate}\paperdate\par}%
  \vspace{10pt}}

% Small headings inside the front matter.
\newcommand{\fronthead}[1]{\par\vspace{6pt}{\sffamily\bfseries\normalsize\color{ink}#1}\par\vspace{-2pt}}
\makeatletter
\newcommand{\contentsblock}{%
  \par\noindent\begin{minipage}{\linewidth}%
  \fronthead{Contents}%
  \setlength{\parskip}{0pt}\setlength{\columnsep}{2.2em}\raggedright
  \begin{multicols}{2}
  \@starttoc{toc}
  \end{multicols}
  \end{minipage}\par}
\makeatother

% ---------- statements ----------
% Theorems, lemmas, propositions and corollaries share one counter.
% Each statement is in a tinted box with an accent rule at the left.
\newtheoremstyle{boxed}{0pt}{0pt}{\normalfont}{0pt}{\sffamily\bfseries\small\color{accent}}{.}{0.6em}{}
\theoremstyle{boxed}
\newtheorem{theorem}{Theorem}[section]
\newtheorem{lemma}[theorem]{Lemma}
\newtheorem{proposition}[theorem]{Proposition}
\newtheorem{corollary}[theorem]{Corollary}

\tcbset{guidebox/.style={enhanced, parbox=false, boxrule=0pt, frame hidden,
  arc=1.5pt, left=11pt, right=10pt, top=7pt, bottom=7pt, before skip=11pt, after skip=9pt}}
% Not breakable: a label at the start of a box that breaks at the top of a page can be lost.
\tcbset{statementbox/.style={guidebox, colback=accent!6,
  borderline west={3pt}{0pt}{accent}}}
\tcolorboxenvironment{theorem}{statementbox}
\tcolorboxenvironment{lemma}{statementbox}
\tcolorboxenvironment{proposition}{statementbox}
\tcolorboxenvironment{corollary}{statementbox}

% The main result of a paper: a framed box with a title bar.
% \begin{maintheorem}{name} uses the counter of the theorems, and \label works after it.
\newtcolorbox{mainresult}[1]{enhanced, parbox=false, colback=white, colframe=accent, boxrule=0.9pt,
  arc=1.5pt, left=11pt, right=11pt, top=8pt, bottom=8pt, title={#1},
  fonttitle=\sffamily\bfseries\small, colbacktitle=accent, coltitle=white,
  before skip=12pt, after skip=12pt}
\newenvironment{maintheorem}[1]%
  {\refstepcounter{theorem}\begin{mainresult}{Theorem~\thetheorem\ (#1)}}%
  {\end{mainresult}}

% ---------- proofs ----------
% A proof has a thin rule at the left, so that the eye finds the statements.
\renewcommand{\proofname}{\upshape\sffamily\bfseries\small\color{slate}Proof}
\renewcommand{\qedsymbol}{\textcolor{slate}{\rule{0.55em}{0.55em}}}
\tcolorboxenvironment{proof}{enhanced, breakable, parbox=false, boxrule=0pt, frame hidden,
  colback=white, borderline west={1.2pt}{0pt}{slate!35}, left=11pt, right=0pt, top=1pt, bottom=2pt,
  before skip=6pt, after skip=11pt}
% A step or a case inside a proof.
\newcommand{\step}[1]{{\sffamily\bfseries\small\color{slate}#1}}

% ---------- idea, notation and fine-print boxes ----------
\newtcolorbox{idea}[1]{guidebox, breakable, colback=warm!6, borderline west={3pt}{0pt}{warm!80},
  before upper={{\sffamily\bfseries\footnotesize\color{warm!75!black}IDEA\enspace\textbar\enspace}{\sffamily\bfseries\small #1}\par}}
\newtcolorbox{notation}[1]{guidebox, breakable, colback=black!4, borderline west={3pt}{0pt}{slate!60},
  fontupper=\small,
  before upper={{\sffamily\bfseries\footnotesize\color{slate}NOTATION\enspace\textbar\enspace}{\sffamily\bfseries\small #1}\par}}
\newtcolorbox{fineprint}{guidebox, breakable, colback=white, borderline west={1.5pt}{0pt}{black!25},
  fontupper=\footnotesize\color{black!78}, left=9pt, top=1pt, bottom=1pt,
  before upper={{\sffamily\bfseries Fine print.}\enspace}}

% ---------- figures ----------
\tikzset{
  vtx/.style={circle, fill=ink, inner sep=1.35pt},
  port/.style={circle, draw=accent, fill=white, line width=0.8pt, inner sep=1.5pt},
  edge/.style={line width=0.55pt, draw=ink},
  hot/.style={line width=1.1pt, draw=warm},
  cool/.style={line width=1.1pt, draw=accent},
  region/.style={draw=accent, fill=accent!7, line width=0.6pt, rounded corners=6pt},
  note/.style={font=\sffamily\footnotesize, text=slate},
}

% ---------- macros of the mathematics ----------
\newcommand{\Z}{\mathcal{Z}}
\newcommand{\PP}{\mathbb{P}}
\newcommand{\EE}{\mathbb{E}}
\newcommand{\comp}{\operatorname{comp}}
\newcommand{\tr}{\operatorname{tr}}
\newcommand{\lgk}[1]{\log^{(#1)}}
\newcommand{\lgs}{\log^{*}}
\newcommand{\sm}{\setminus}
\newcommand{\half}{\tfrac12}
\newcommand{\cO}{\mathcal{O}}


\renewcommand{\papertitle}{The forest proportion\\ of a random even subgraph}
\renewcommand{\papersubtitle}{A second-moment estimate for the edges away from a small marked set}
\renewcommand{\paperdate}{Draft of 12 November 2026}
\renewcommand{\paperrunning}{The forest proportion of a random even subgraph}
\hypersetup{pdftitle={The forest proportion of a random even subgraph},pdfauthor={KNT, Astra, Fable 5.2}}

\begin{document}
\maketitleblock

An even subgraph of a graph is a set of edges that meets every vertex an even number of times.
Let $\Gamma$ be a connected graph of large rank $r$ in which a small set $P$ of vertices is
marked and all other vertices have degree three. We prove that, for a uniformly random even
subgraph of $\Gamma$, the edges without an end in $P$ form a forest with probability at most
\[
  \frac12+\frac{20}{\log\log\log r}.
\]
The constant $\frac12$ cannot be lowered. Two companion papers use this \emph{forest estimate}.
The first determines the minimum number of edges of a pancyclic graph up to an additive
constant~\cite{Pan}, and the second determines the maximum number of cycles of a graph of large
rank~\cite{Max}. Astra and Fable~5.2 are language models. Section~\ref{sec:status} states the
part of each author. The estimate has a formal proof in Lean~4, and the Lean kernel has checked
it.

\fronthead{The proof in five sentences}
\begin{enumerate}[leftmargin=1.6em]
\item A cycle $C$ away from $P$ is a component of the random even subgraph $\xi$ with
probability at least $2^{-|C|}$, and then the edges of $\xi$ away from $P$ do not form a forest.
\item For two disjoint cycles, these two events are independent if the cycles have one
\emph{link}, and their joint probability is twice the product if they have two links; a cycle
has few disjoint partners with three or more links.
\item If $\Gamma-P$ has many pairwise disjoint short cycles, the second moment of the number of
these events gives the estimate. So we may delete a small maximal family of short cycles, and a
graph $J$ of girth more than $\frac1{10}\log r$ remains.
\item A count of closed walks without backtracking, by the largest eigenvalue of the
non-backtracking matrix, shows that the cycles of $J$ with between $6\log r$ and
$\log r\sqrt{\log\log r}$ edges have a total weight $\mu=\sum_C2^{-|C|}$ of order
$\log\log\log r$.
\item By the second sentence, a weighted number $Z$ of the events for these cycles has
$\EE Z^{2}\le2\mu^{2}+2\mu$, and the Cauchy--Schwarz inequality gives
$\PP(Z=0)\le\half+1/(2\mu)$.
\end{enumerate}

\contentsblock

\section{The estimate}\label{sec:intro}

\subsection{The result}

Throughout, $\log$ is the binary logarithm, and $\lgk{j}$ is its $j$-fold iterate. A graph may
have loops and parallel edges. We say \emph{multigraph} to stress this, and \emph{simple graph}
to exclude it. A loop adds two to the degree of its vertex. A multigraph is \emph{subcubic} if
every vertex has degree at most three.

For a multigraph $G$, an \emph{even subgraph} is a set of edges that meets every vertex an even
number of times, where a loop meets its vertex twice. With symmetric difference as addition, the
even subgraphs form a vector space $\Z(G)$ over $\mathbb{F}_2$. Its dimension is
\[
  \beta(G)=|E(G)|-|V(G)|+\comp(G),
\]
where $\comp(G)$ is the number of components. We call $\beta(G)$ the \emph{rank} of $G$; other
names are cyclomatic number and cycle rank. A \emph{cycle} is a connected $2$-regular subgraph.
A loop is a cycle of length one, and two parallel edges form a cycle of length two. We identify a
subgraph without isolated vertices with its edge set. Then a cycle is a nonzero element of
$\Z(G)$, and a \emph{forest} is a set of edges that contains no cycle.

Let $P$ be a set of vertices of $G$. We write $G[P]$ for the induced subgraph and $G-P$ for
$G[V(G)\sm P]$. For an even subgraph $\xi$, the \emph{outside part} of $\xi$ is
$\xi\cap E(G-P)$, the set of its edges without an end in $P$ (Figure~\ref{fig:outside}). We
write $q(G,P)$ for the proportion of the even subgraphs of $G$ whose outside part is a forest.
So $q(G,P)$ is the probability that the outside part of $\xi$ is a forest, when $\xi$ is a
uniformly random element of $\Z(G)$.

\begin{figure}[ht]
\centering
\begin{tikzpicture}[x=0.82cm,y=0.82cm]
  \foreach \sh/\which in {0/0,8.6/1}{
  \begin{scope}[xshift=\sh*0.82cm]
    \foreach \i in {0,...,5} \coordinate (w\i) at ({60*\i+30}:0.9);
    \coordinate (a) at (1.9,0.75);
    \coordinate (b) at (3.0,0);
    \coordinate (c) at (1.9,-0.75);
    \coordinate (d) at (-1.9,0.75);
    \coordinate (e) at (-3.0,0);
    \coordinate (f) at (-1.9,-0.75);
    % all edges, thin
    \draw[edge] (w0)--(w1)--(w2)--(w3)--(w4)--(w5)--cycle;
    \draw[edge] (w0)--(a) (w5)--(c) (w2)--(d) (w3)--(f);
    \draw[edge] (a)--(b)--(c)--(a) (d)--(e)--(f)--(d);
    \draw[edge] (b) .. controls (2.9,1.7) and (1.0,1.5) .. (w1);
    \draw[edge] (e) .. controls (-2.9,-1.7) and (-1.0,-1.5) .. (w4);
    \ifnum\which=0
      % a cycle through P
      \draw[hot] (a)--(b)--(c);
      \draw[white, line width=1.9pt] (w5)--(w0) (w0)--(a) (c)--(w5);
      \draw[hot, dash pattern=on 3pt off 2pt] (w5)--(w0) (w0)--(a) (c)--(w5);
    \else
      % the hexagon and a triangle
      \draw[hot] (a)--(b)--(c)--cycle;
      \draw[white, line width=1.9pt] (w0)--(w1)--(w2)--(w3)--(w4)--(w5)--cycle;
      \draw[hot, dash pattern=on 3pt off 2pt] (w0)--(w1)--(w2)--(w3)--(w4)--(w5)--cycle;
    \fi
    \foreach \p in {w0,w1,w2,w3,w4,w5,a,b,c,d,e,f} \node[vtx] at (\p) {};
    \node[note, text=accent] at (0,0) {$P$};
  \end{scope}}
  \node[note, text=warm] at (0,-2.05) {the outside part is a path};
  \node[note, text=warm] at (8.6,-2.05) {the outside part is a triangle};
\end{tikzpicture}
\caption{A cubic graph, a set $P$ of six vertices, and two even subgraphs, drawn thick. The
edges with an end in $P$ are dashed. Left: a cycle that meets $P$. Its outside part is a path,
so it is a forest. Right: the union of two cycles. Its outside part is a triangle, so it is not
a forest.}
\label{fig:outside}
\end{figure}

\begin{maintheorem}{the forest estimate}\label{thm:forest}
There is $r_0$ such that the following holds for all $r\ge r_0$. Let $\Gamma$ be a connected
subcubic multigraph of rank $r$. Let $P$ be a set of at most $2r/\log^{2}r$ vertices such that
every vertex outside $P$ has degree three and $\Gamma[P]$ has at most $(\log r)/100$ components.
Then
\[
  q(\Gamma,P)\;\le\;\frac12+\frac{20}{\lgk{3}r}.
\]
\end{maintheorem}

The estimate and its proof are from~\cite{KNT}. The proof gives no value for $r_0$.

\begin{fineprint}
The constant $\half$ cannot be lowered. Let $k\ge2$, and take the vertices $a_i,b_i,c_i,d_i$
for $i\in\mathbb{Z}/k$ with the edges $a_ic_i$, $a_id_i$, $b_ic_i$, $b_id_i$, $c_ia_{i+1}$ and
$d_ib_{i+1}$. This is a ring of $k$ copies of $K_{2,2}$. It is cubic, and its rank is
$r=2k+1$. Consider the cycles that contain exactly one of the edges $c_ia_{i+1}$, $d_ib_{i+1}$
for every $i$. There are $2^{k}$ choices for these edges, and then two paths through each copy
of $K_{2,2}$. So there are $2^{2k}=2^{r-1}$ of these cycles, and each of them meets the set
$P=\{a_0,b_0,c_0,d_0\}$. The outside part of a cycle that meets $P$ is a proper subset of a
cycle, hence a forest. So $q(\Gamma,P)\ge\half$ for this graph $\Gamma$ of rank $r$.
\end{fineprint}

The second form of the estimate has no condition on the degrees, and a subgraph is marked in
the place of a set of vertices. Let $W$ be a subgraph of a graph $G$. For an even subgraph
$\xi$ of $G$ we call $\xi\sm E(W)$ its \emph{outside part with respect to $W$}. Let $p(G,W)$ be
the probability that $\xi\sm E(W)$ is a forest, for a uniformly random $\xi\in\Z(G)$.

\begin{corollary}[Simple graphs]\label{cor:simple}
There is $r_0$ such that the following holds for all $r\ge r_0$. Let $G$ be a connected simple
graph of rank $r$. Let $W\subseteq G$ be a nonempty union of cycles with at most $r/\log^{2}r$
edges and at most $(\log r)/100$ components. Then
\[
  p(G,W)\;\le\;\frac12+\frac{20}{\lgk{3}r}.
\]
\end{corollary}

Sections~\ref{sec:tools} to~\ref{sec:proof} prove Theorem~\ref{thm:forest}, and
Section~\ref{sec:simple} derives the corollary.

\subsection{Outline of the proof}

\begin{idea}{the proof in three steps}
\begin{enumerate}[leftmargin=1.5em]
\item If $\Gamma-P$ has many pairwise disjoint short cycles, the estimate follows from
Lemma~\ref{lem:smallcuts}. So we may delete a small maximal family of them, and a graph $J$ of
large girth remains.
\item The graph $J$ has many cycles with between $6\log r$ and $\log r\sqrt{\log\log r}$ edges.
Their total weight $\sum_C2^{-|C|}$ is at least $(\frac18-o(1))\ln\log\log r$
(Lemma~\ref{lem:weight}).
\item For such a cycle $C$, the event that $C$ is a component of $\xi$ has probability at least
$2^{-|C|}$, and then the outside part of $\xi$ is not a forest. Most pairs of these events are
correlated by a factor of at most two (Lemma~\ref{lem:link}).
\end{enumerate}
With the factor two of step~3, the second moment of a weighted number $Z$ of these events is at
most $2\mu^{2}+K\mu$, and Lemma~\ref{lem:second} then gives $\PP(Z=0)\le\half+K/(4\mu)$.
\end{idea}

The three hypotheses on $P$ enter as follows. The vertices outside $P$ have degree three. So a
cycle of $\Gamma-P$ that is contained in $\xi$ is a component of $\xi$, and $J$ has about
$2^{\ell}$ closed walks without backtracking of length $\ell$. The bound on $|P|$ keeps the
number of vertices of $J$ of degree less than three small against $r/\log r$, which
Lemma~\ref{lem:weight} needs for cycles of length about $\log r$. Each component of $\Gamma[P]$
can give one more link between two cycles, and the bound on the number of components keeps the
number of links below $\frac1{90}\log r$ (Step~2 in the proof of
Proposition~\ref{prop:loggirth}).

\subsection{Two applications}\label{sec:applications}

Two companion papers use the estimate. We state their results without proof, and we state how
the hypotheses of the estimate are satisfied.

\paragraph{Pancyclic graphs~\cite{Pan}.}
A simple graph on $n$ vertices is \emph{pancyclic} if it has a cycle of every length from $3$
to $n$. Let $n+h(n)$ be the least number of edges of such a graph. Then
$h(n)=\log n+\lgs n+O(1)$, where $\lgs$ is the iterated logarithm. This answers Problem~\#1016
of Erd\H{o}s~\cite{Bloom}. The lower bound uses Corollary~\ref{cor:simple}. Let $G$ be a graph
of rank $r\ge2^{256R}$ with a cycle of every length from $3$ to a large number. Cycles of the
first lengths form a subgraph $W$ that is a union of at most $2R$ cycles and has at most
$2^{2R+1}$ edges. So $W$ has at most $\sqrt r\le r/\log^{2}r$ edges and at most
$2R\le(\log r)/100$ components. These cycles can be disjoint, and for this reason the estimate
allows several components. A cycle of $G$ with an edge in $W$ has a forest as its outside part,
and the cycles with a given outside part of this kind have at most $|E(W)|$ different lengths.
By the corollary, at most $\half+o(1)$ of the outside parts are forests. So the largest possible
density of cycle lengths decreases by a factor of almost two from the rank $R$ to the rank
$2^{256R}$. This gives the term $\lgs n$.

\paragraph{The maximum number of cycles~\cite{Max}.}
Let $\psi(r)$ be the largest number of cycles of a graph of rank $r$. Entringer and
Slater~\cite{ES} conjectured that $\psi(r)\sim2^{r-1}$. In fact $\psi(r)=H(r)$ for all large
$r$, where $H(r)=2^{r-1}+(1+o(1))\,r\,2^{\lfloor(r-1)/2\rfloor}$ is the number of cycles of
cubic graphs of Horak and Kim~\cite{HK}. The proof uses Theorem~\ref{thm:forest} only when
$\Gamma[P]$ is connected. Some connected cubic multigraph $G$ of rank $r$ has $\psi(r)$ cycles.
Let $P$ be a set of $\lceil4\log r\rceil$ vertices such that $G[P]$ is connected. The outside
part of a cycle that meets $P$ is a proper subset of a cycle, hence a forest. So at most
$q(G,P)\,2^{r}$ cycles meet $P$. If $G$ has no bridge, at most $2^{r-1-|P|/2}$ cycles avoid
$P$. This gives $\psi(r)\le(\half+o(1))2^{r}$. For the exact value, \cite{Max} applies the
theorem also to parts of $G$ with a small boundary. Where a bound below $\half$ is necessary,
it uses Lemmas~\ref{lem:link} to~\ref{lem:smallcuts} and Proposition~\ref{prop:loggirth}
directly. For this purpose Proposition~\ref{prop:loggirth} has the hypothesis
$q(\Gamma,P)\ge\frac25$, which is weaker than $q(\Gamma,P)>\half$.

\subsection{Authorship and verification}\label{sec:status}

KNT is a person. Astra and Fable~5.2 are the language models GPT-6 Astra and Claude Fable~5.2.
KNT directed the work at each step. KNT obtained the forest estimate and its proof in September
2026 with Astra, for the problem on pancyclic graphs~\cite{KNT}. Fable~5.2 wrote this text.

Theorem~\ref{thm:forest} and Corollary~\ref{cor:simple} have formal proofs in Lean~4. The
statements are \lean{Erdos1016.ShortProof.eventually_marked_forest_bound_logSquared} and
\lean{Erdos1016.ShortProof.fewComponentForestEstimate_logSquared}. In the first statement the
rank is $2^{x}$, and the bound is $\half+20/\log\log x$ for all sufficiently large real $x$.
The Lean kernel accepts the proofs with the axioms \lean{propext}, \lean{Classical.choice} and
\lean{Quot.sound} only. The files are in~\cite{Repo}, at commit \texttt{3a80aa9}.

The formal proof has the steps of this paper. In one place its method is different: for the
links of two cycles, which is Step~2 in the proof of Proposition~\ref{prop:loggirth}, it uses
Lemma~\ref{lem:smallcuts}, and so it supposes that $q(\Gamma,P)>\half$. The kernel check covers
the two formal statements, and it does not cover the text of this paper. No referee has read
the text.

\section{Cut events and links}\label{sec:tools}

\begin{notation}{regions and the random even subgraph}
\begin{itemize}[leftmargin=1.2em]
\item For a set $A$ of vertices of a multigraph $G$, $\delta A$ is the set of the edges with
exactly one end in $A$.
\item A \emph{region} is a nonempty set $A$ of vertices such that $G[A]$ is connected. A region
is \emph{cyclic} if $G[A]$ contains a cycle.
\item Throughout, $\xi$ is a uniformly random element of $\Z(G)$.
\item For a region $A$, the \emph{cut event} $\cO_A$ is the event $\xi\cap\delta A=\emptyset$.
\end{itemize}
\end{notation}

\begin{lemma}[Second moment]\label{lem:second}
Let $Z\ge0$ be a random variable with $\mu=\EE Z>0$ and $\EE Z^{2}\le2\mu^{2}+K\mu$. Then
$\PP(Z=0)\le\half+K/(4\mu)$.
\end{lemma}

\begin{proof}
By the Cauchy--Schwarz inequality $\mu^{2}\le\EE Z^{2}\cdot\PP(Z>0)$, so
$\PP(Z=0)\le1-\mu/(2\mu+K)\le\half+K/(4\mu)$.
\end{proof}

Let $A$ and $B$ be disjoint regions of a connected multigraph $G$. A \emph{link} between them is an
edge with one end in $A$ and the other in $B$, or a component of $G-(A\cup B)$ that is joined by an
edge to $A$ and by an edge to $B$ (Figure~\ref{fig:links}). Let $\ell(A,B)$ be the number of links.

\begin{figure}[ht]
\centering
\begin{tikzpicture}
  \draw[region] (0,0) ellipse [x radius=0.75, y radius=1.0];
  \draw[region] (6,0) ellipse [x radius=0.75, y radius=1.0];
  \draw[draw=warm, fill=warm!8, line width=0.6pt] (3,-0.7) ellipse [x radius=0.85, y radius=0.38];
  \draw[draw=slate!70, fill=black!4, line width=0.6pt] (-2.9,0.2) ellipse [x radius=0.85, y radius=0.38];
  \draw[draw=slate!70, fill=black!4, line width=0.6pt] (8.9,-0.2) ellipse [x radius=0.85, y radius=0.38];
  \node[text=accent] at (0,0) {$A$};
  \node[text=accent] at (6,0) {$B$};
  \node[text=warm] at (3,-0.7) {$Y_1$};
  \node[text=slate] at (-2.9,0.2) {$Y_2$};
  \node[text=slate] at (8.9,-0.2) {$Y_3$};
  \draw[hot] (0.45,0.55) -- node[above=-1pt, text=warm] {\small$e$} (5.55,0.55);
  \draw[edge] (0.5,-0.45) -- (2.25,-0.58);
  \draw[edge] (0.4,-0.75) -- (2.3,-0.86);
  \draw[edge] (5.5,-0.5) -- (3.75,-0.63);
  \draw[edge] (-0.5,0.35) -- (-2.12,0.32);
  \draw[edge] (-0.55,0.05) -- (-2.1,0.1);
  \draw[edge] (6.55,-0.15) -- (8.07,-0.2);
  \foreach \p in {(0.45,0.55),(5.55,0.55),(0.5,-0.45),(2.25,-0.58),(0.4,-0.75),(2.3,-0.86),(5.5,-0.5),
                  (3.75,-0.63),(-0.5,0.35),(-2.12,0.32),(-0.55,0.05),(-2.1,0.1),(6.55,-0.15),(8.07,-0.2)}
    \node[vtx] at \p {};
  \node[note, text=warm] at (3,1.05) {two links: the edge $e$ and the component $Y_1$};
  \node[note] at (-2.9,-0.45) {not a link};
  \node[note] at (8.9,-0.85) {not a link};
\end{tikzpicture}
\caption{Links between two regions $A$ and $B$. The edge $e$ and the component $Y_1$ of
$G-(A\cup B)$ are links. The components $Y_2$ and $Y_3$ are joined to one region only, so they are
not links. Here $\ell(A,B)=2$.}
\label{fig:links}
\end{figure}

\begin{lemma}[Link formula]\label{lem:link}
Let $A$ and $B$ be disjoint regions of a connected multigraph $G$. Then
$\PP(\cO_A)=2^{\comp(G-A)-|\delta A|}\ge2^{-|\delta A|}$ and
\[
  \PP(\cO_A\cap\cO_B)=2^{\ell(A,B)-1}\,\PP(\cO_A)\,\PP(\cO_B).
\]
Conditionally on $\cO_A$, the set $\xi\cap E(G[A])$ is uniform on $\Z(G[A])$. Conditionally on
$\cO_A\cap\cO_B$, the sets $\xi\cap E(G[A])$ and $\xi\cap E(G[B])$ are independent and uniform on
$\Z(G[A])$ and $\Z(G[B])$.
\end{lemma}

\begin{proof}
The even subgraphs that avoid $\delta A$ are the sums of an even subgraph of $G[A]$ and an even
subgraph of $G-A$. So $\PP(\cO_A)=2^{\beta(G[A])+\beta(G-A)-\beta(G)}$, and a count of edges and
vertices shows that the exponent is $\comp(G-A)-|\delta A|$. In the same way
\[
  \PP(\cO_A\cap\cO_B)=2^{\beta(G[A])+\beta(G[B])+\beta(G-A-B)-\beta(G)}
     =2^{\comp(G-A-B)+1-|\delta A|-|\delta B|+e(A,B)},
\]
where $e(A,B)$ is the number of edges between $A$ and $B$. Let $a$, $b$ and $k$ be the numbers of
components of $G-A-B$ that are joined to $A$ only, to $B$ only, and to both. Then
$\comp(G-A-B)=a+b+k$, $\comp(G-A)=a+1$ and $\comp(G-B)=b+1$. The formula follows, because
$\ell(A,B)=k+e(A,B)$. The conditional laws are read off from the two decompositions as direct sums.
\end{proof}

\begin{idea}{the correlation of two cut events}
By Lemma~\ref{lem:link}, two cut events are independent when the regions have one link, and
their joint probability is twice the product when the regions have two links. The constant
$\half$ of the forest estimate comes from this factor two, through Lemma~\ref{lem:second}.
\end{idea}

Pairs with three or more links are rare:

\begin{lemma}[Exceptional partners]\label{lem:exc}
Let $A$ be a region of a connected multigraph $G$, and let $B_1,\dots,B_k$ be pairwise disjoint
regions, disjoint from $A$, with $\ell(A,B_i)\ge3$ for all $i$. Then $k\le|\delta A|$.
\end{lemma}

\begin{proof}
Contract $A$ to a vertex $a$ and each $B_i$ to a vertex $b_i$, and delete the loops that arise.
This changes neither $d=|\delta A|$ nor the numbers $\ell(A,B_i)$, because the sets $B_j$ with
$j\ne i$ are connected and are contracted inside components of $G-(A\cup B_i)$. Let $F$ be a
spanning forest of $G-a$, and for every edge $e=au$ at $a$ attach a new leaf $\tau_e$ to $u$. Call
the resulting forest $F^{+}$ and its $d$ new leaves \emph{terminals}.

Fix $i$. Every link between $a$ and $b_i$ yields a path of $F^{+}$ from $b_i$ to a terminal. For an
edge $e=ab_i$ it is the edge $b_i\tau_e$. For a component $Y$ of $G-\{a,b_i\}$ joined to both,
choose an edge $e=au$ with $u\in Y$ and take the path of $F^{+}$ from $b_i$ to $\tau_e$; after its
first vertex it stays in $Y\cup\{\tau_e\}$. Paths that come from different links start with
different edges at $b_i$, and the union of any two of them is a path between two terminals. Let
$F^{\circ}\subseteq F^{+}$ be the union of all paths between terminals. As $\ell(a,b_i)\ge3$, the
vertex $b_i$ has degree at least $3$ in $F^{\circ}$. All leaves of $F^{\circ}$ are terminals, and a
forest with at most $d$ leaves has at most $d-2$ vertices of degree at least $3$. Hence $k\le d-2$
if $k\ge1$.
\end{proof}

\begin{lemma}[Many small cuts]\label{lem:smallcuts}
Let $A_1,\dots,A_k$ be pairwise disjoint cyclic regions of a connected multigraph $G$ with
$|\delta A_i|\le d$ for all $i$. Then
\[
  \PP\bigl(\xi\cap E(G[A_i])\text{ is a forest for every }i\bigr)\;\le\;\frac12+\frac{(d+1)\,2^{d-1}}{k}.
\]
\end{lemma}

\begin{proof}
Let $E_i$ be the event that $\cO_{A_i}$ holds and $\xi\cap E(G[A_i])\ne\emptyset$. On $E_i$ the set
$\xi\cap E(G[A_i])$ is a nonempty even subgraph, hence not a forest. So it suffices to bound
$\PP(Z=0)$ for $Z=\sum_i\mathbf{1}_{E_i}$. By Lemma~\ref{lem:link},
$p_i:=\PP(E_i)=\PP(\cO_{A_i})\bigl(1-2^{-\beta(G[A_i])}\bigr)\ge2^{-d-1}$ and
$\PP(E_i\cap E_j)=2^{\ell(A_i,A_j)-1}p_ip_j$ for $i\ne j$. If $\ell(A_i,A_j)\le2$, this is at most
$2p_ip_j$. Otherwise we use $\PP(E_i\cap E_j)\le p_i$; by Lemma~\ref{lem:exc} this happens for at
most $d$ indices $j$, for each $i$. Hence $\EE Z^{2}\le\mu+2\mu^{2}+d\mu$ with
$\mu=\EE Z\ge k\,2^{-d-1}$, and Lemma~\ref{lem:second} gives the claim.
\end{proof}

\section{Cycles of moderate length in a graph of large girth}\label{sec:moderate}

\begin{idea}{the content of this section}
Delete few vertices from $\Gamma$ such that a simple graph $J$ of girth more than $D$ remains.
Almost all vertices of $J$ have degree three. So $J$ has about $2^{\ell}$ closed walks without
backtracking of each even length $\ell$, and for moderate $\ell$ most of these walks go once
around a cycle. Hence the cycles $C$ of $J$ with $\ell_0\le|C|\le L$ have a total weight
$\sum_C2^{-|C|}$ of about $\sum_\ell1/(2\ell)$, which is about $\frac14\ln(L/\ell_0)$
(Lemma~\ref{lem:weight}). Proposition~\ref{prop:loggirth} then bounds the second moment of the
number of these cycles that are components of $\xi$, counted with weights.
\end{idea}

\begin{notation}{the setting of this section}
\begin{itemize}[leftmargin=1.2em]
\item $\Gamma$ is a connected subcubic multigraph, and $U$ is a set of its vertices such that
every vertex outside $U$ has degree $3$.
\item $J=\Gamma-U$ is a simple graph of girth greater than $D$, where $D\ge3$ is an integer, and
$s=\lfloor(D-1)/2\rfloor$. The \emph{girth} is the length of a shortest cycle.
\item $L\ge1$ is an integer, and $\mathcal{C}$ is the set of the cycles of $J$ of length at most
$L$.
\item $v$ is the number of vertices of $J$, and $d_J=\sum_{u\in V(J)}(3-\deg_Ju)$ is its
\emph{deficit}. This is the number of edges between $V(J)$ and $U$, so $d_J\le3|U|$.
\end{itemize}
\end{notation}

\subsection{Walks without backtracking}

A walk $u_0u_1\cdots u_\ell$ is \emph{non-backtracking} if $u_{i+1}\ne u_{i-1}$ for $0<i<\ell$. Let
$N_\ell(u,w)$ be the number of non-backtracking walks of length $\ell$ from $u$ to $w$ in $J$.

\begin{lemma}[Walks without backtracking]\label{lem:walks}
If $\ell>s$, then $N_\ell(u,w)\le3\cdot2^{\ell-s-1}$ for all $u,w$. Consequently, for every vertex
$w$ of $J$,
\begin{equation}\label{eq:load}
  \sum_{C\in\mathcal{C},\ C\ni w}2^{-|C|}\;\le\;L\,2^{-s}.
\end{equation}
\end{lemma}

\begin{proof}
There are at most $3\cdot2^{\ell-s-1}$ ways to choose the first $\ell-s$ steps. The last $s$ steps
form a non-backtracking walk of length $s$ between two given vertices. Such a walk is a path,
because $s$ is less than the girth. There is at most one such path, because two of them would give
a cycle of length at most $2s<D$. For the second statement, a cycle of length $\ell$ through $w$
can be traversed from $w$ in two directions. So there are at most $\half N_\ell(w,w)$ of them, and
$\ell>D>s$. Hence the cycles of each length $\ell\le L$ contribute at most $\frac34\,2^{-s}$.
\end{proof}

Let $B$ be the non-backtracking matrix of $J$. Its rows and columns are indexed by the directed
edges $(i\to j)$ of $J$, two for each edge. The entry in row $(i\to j)$ and column $(j'\to k)$ is
$1$ if $j'=j$ and $k\ne i$, and it is $0$ otherwise.

\begin{lemma}[Closed walks]\label{lem:spectrum}
Suppose that $d_J<v$. Then, for every even $\ell\ge2$,
\[
  \tr B^{\ell}\;\ge\;\bigl(2-d_J/v\bigr)^{\ell}-3v\,2^{\ell/2}.
\]
\end{lemma}

\begin{proof}
Let $A$ be the adjacency matrix of $J$ and $\Delta$ the diagonal matrix of its degrees, and put
$M(\lambda)=\lambda^{2}I-\lambda A+\Delta-I$.

(1) \emph{If $M(\lambda)g=0$ for a real $\lambda>1$ and a vector $g\ne0$, then $\lambda$ is an
eigenvalue of $B$.} A direct computation shows that $f(i\to j)=\lambda g(j)-g(i)$ satisfies
$Bf=\lambda f$. And $f=0$ would give $g(i)=\lambda g(j)=\lambda^{2}g(i)$ along every edge, hence
$g=0$; at an isolated vertex, the equation $M(\lambda)g=0$ is $(\lambda^{2}-1)g=0$.

(2) \emph{If $Bf=\lambda f$ with $f\ne0$ and $\lambda^{2}\ne1$, then $M(\lambda)g=0$ for
$g(j)=\sum_{k\sim j}f(j\to k)$, and $g\ne0$.} The eigenvalue equation at $(i\to j)$ and at
$(j\to i)$ reads $\lambda f(i\to j)=g(j)-f(j\to i)$ and $\lambda f(j\to i)=g(i)-f(i\to j)$.
Eliminating $f(j\to i)$ gives $(\lambda^{2}-1)f(i\to j)=\lambda g(j)-g(i)$. This shows that
$g\ne0$. Summing it over the neighbours $j$ of $i$ gives $M(\lambda)g=0$.

(3) \emph{$B$ has a real eigenvalue $\lambda_1\ge2-d_J/v$.} The matrix
$M(2)=2(\Delta-A)+(3I-\Delta)$ is positive semidefinite, while
$\mathbf{1}^{\!\top}M(\lambda)\mathbf{1}=(\lambda-1)(v\lambda-2v+d_J)$ vanishes at
$\lambda^{*}=2-d_J/v>1$. The smallest eigenvalue of the symmetric matrix $M(\lambda)$ is continuous
in $\lambda$, at most $0$ at $\lambda^{*}$ and at least $0$ at $2$. So $M(\lambda_1)$ is singular for
some $\lambda_1\in[\lambda^{*},2]$, and (1) applies.

(4) \emph{Every non-real eigenvalue $\lambda$ of $B$ satisfies $|\lambda|\le\sqrt2$.} With $g$ as
in~(2), $g^{*}M(\lambda)g=0$ is a quadratic equation $a\lambda^{2}-b\lambda+c'=0$ with real
coefficients $a=g^{*}g>0$, $b=g^{*}Ag$ and $c'=g^{*}(\Delta-I)g\le2a$. Its roots are $\lambda$ and
$\bar\lambda$, so $|\lambda|^{2}=c'/a\le2$.

Now $\tr B^{\ell}$ is the sum of the $\ell$th powers of the $2|E(J)|\le3v$ eigenvalues of $B$. For
even $\ell$ the real eigenvalues contribute nonnegative terms, one of which is at least
$(2-d_J/v)^{\ell}$, and the non-real ones contribute at least $-3v\,2^{\ell/2}$ in total.
\end{proof}

Facts (2) and (4) are part of the Ihara--Bass theory~\cite{Bass,KS}. The test vector $\mathbf{1}$
in~(3) is from~\cite{KNT}. With this test vector the argument needs no assumption of regularity
or expansion.

\begin{lemma}[Cycle weight]\label{lem:weight}
Suppose that $d_J<v$, and let $\ell_0\le L$ be a positive integer. The total weight
$\mu:=\sum_{C\in\mathcal{C}}2^{-|C|}$ satisfies
\[
  \Bigl(1+\tfrac32\,L^{2}\,2^{-s}\Bigr)\,\mu\;\ge\;
  \Bigl(1-\frac{L\,d_J}{2v}-3v\,2^{-\ell_0/2}\Bigr)\sum_{\ell}\frac1{2\ell},
\]
where the sum is over the even $\ell$ with $\ell_0\le\ell\le L$.
\end{lemma}

\begin{proof}
$\tr B^{\ell}$ is the number of closed walks $u_0u_1\cdots u_\ell=u_0$ in $J$ with
$u_{i+1}\ne u_{i-1}$ for all $i$ modulo $\ell$. A cycle of length $\ell$ gives $2\ell$ of them; let
$n_\ell$ be the number of the others. Then
\[
  T:=\sum_{\ell\le L}\frac{\tr B^{\ell}}{2\ell\,2^{\ell}}=\mu+\sum_{\ell\le L}\frac{n_\ell}{2\ell\,2^{\ell}} .
\]
Consider a walk counted by $n_\ell$, and let $j$ be the least index with
$u_j\in\{u_0,\dots,u_{j-1}\}$, say $u_j=u_i$. Then $j<\ell$, the walk $C=u_iu_{i+1}\cdots u_j$ is a
cycle, of length $q$ say, and $u_ju_{j+1}\cdots u_\ell u_1\cdots u_i$ is a non-backtracking closed
walk of length $\ell-q>0$ from $u_i$ to itself. It contains a cycle, so $\ell-q>D>s$. The original
walk is determined by $C$, by the vertex $u_i$ and a direction on $C$ ($2q$ choices), by the
complementary walk (at most $3\cdot2^{\ell-q-s-1}$ choices, by Lemma~\ref{lem:walks}), and by the
position of $u_0$ on the latter (at most $\ell$ choices). Hence
\[
  \sum_{\ell\le L}\frac{n_\ell}{2\ell\,2^{\ell}}
  \;\le\;\sum_{\ell\le L}\ \sum_{C\in\mathcal{C}}\tfrac32\,|C|\,2^{-|C|-s}
  \;\le\;\tfrac32\,L^{2}\,2^{-s}\,\mu ,
\]
and $T\le(1+\frac32L^{2}2^{-s})\mu$. On the other hand, all terms of $T$ are nonnegative. For even
$\ell$ with $\ell_0\le\ell\le L$, Lemma~\ref{lem:spectrum} gives
\[
  2^{-\ell}\tr B^{\ell}\;\ge\;\Bigl(1-\frac{d_J}{2v}\Bigr)^{\ell}-3v\,2^{-\ell/2}
  \;\ge\;1-\frac{L\,d_J}{2v}-3v\,2^{-\ell_0/2}.\qedhere
\]
\end{proof}

\subsection{The events and their second moment}\label{sec:events}

Let $\xi$ be uniform on $\Z(\Gamma)$. For $C\in\mathcal{C}$ let $I_C$ be the event that $C$ is a
component of $\xi$. As the vertices of $C$ have degree $3$ in $\Gamma$, this is the event
$E(C)\subseteq\xi$. It is also the event that $\cO_{V(C)}$ holds and
$\xi\cap E(\Gamma[V(C)])=E(C)$. Put
\[
  a_C=2^{-|C|},\qquad p_C=\PP(I_C),\qquad X_C=\frac{a_C}{p_C}\,\mathbf{1}_{I_C},\qquad
  Z=\sum_{C\in\mathcal{C}}X_C .
\]
Because of the factor $a_C/p_C$, the mean of $X_C$ is the combinatorial weight $a_C$ for every
value of $p_C$. The event $I_C$ is given by $|C|$ linear conditions on $\xi\in\Z(\Gamma)$,
which are satisfied by $\xi=C$. So $p_C\ge a_C$. Hence $0\le X_C\le1$, and $\EE Z=\mu$. If
$Z\ne0$, then $\xi$ contains a cycle of $J$.

Let $C\ne C'$ be in $\mathcal{C}$. If they share a vertex, then $I_C\cap I_{C'}=\emptyset$. If they
are disjoint, Lemma~\ref{lem:link}, applied to the regions $V(C)$ and $V(C')$ of $\Gamma$, gives
$\PP(I_C\cap I_{C'})=2^{\ell(C,C')-1}p_Cp_{C'}$ and therefore
\begin{equation}\label{eq:pair}
  \EE(X_CX_{C'})=a_Ca_{C'}\,2^{\ell(C,C')-1},
\end{equation}
where $\ell(C,C')$ is the number of links between $V(C)$ and $V(C')$ in $\Gamma$. Let $W_2$ be
the sum of $a_Ca_{C'}$ over the ordered pairs of disjoint cycles of $\mathcal{C}$ with exactly two
links. Then $W_2\le\mu^{2}$.

\subsection{Graphs of logarithmic girth}

\begin{proposition}[Logarithmic girth]\label{prop:loggirth}
There is $r_2$ such that the following holds for all $r\ge r_2$, with $x=\log r$, $y=\log x$ and
$\eta=r^{-1/60}$. Let $\Gamma$ be a connected subcubic multigraph of rank $r$, and let $P$ be a
set of vertices as in Theorem~\ref{thm:forest}, with $q(\Gamma,P)\ge\frac25$. Let $D$ and $L$ be
integers with $x/16\le D\le x/8$ and $7x\le L\le x\sqrt y$. Let $S$ be the union of the vertex
sets of a maximal family of pairwise disjoint cycles of length at most $D$ in $\Gamma-P$, and
suppose that $|S|\le r/x^{2}$. Then $J=\Gamma-(P\cup S)$ is a simple graph of girth greater than
$D$. So $\Gamma$ and $U=P\cup S$ are as in this section, and
\begin{enumerate}
\item[\textup{(i)}] $L^{2}\,2^{-s}\le\eta$, so that the left side of~\eqref{eq:load} is at most
$\eta/L$ for every vertex $w$ of $J$;
\item[\textup{(ii)}] $\mu\ge\frac14\ln\dfrac{L}{6x}-\dfrac1{\sqrt x}$;
\item[\textup{(iii)}] $\EE Z^{2}\le\mu^{2}+W_2+(1+\eta)\,\mu$.
\end{enumerate}
\end{proposition}

\begin{figure}[ht]
\centering
\begin{tikzpicture}[x=1cm,y=1cm]
  \draw[cool] (0,0) ellipse [x radius=0.7, y radius=1.4];
  \draw[cool] (6.4,0) ellipse [x radius=0.7, y radius=1.4];
  \node[note, text=accent] at (-1.05,0) {$C$};
  \node[note, text=accent] at (7.5,0) {$C'$};
  % the ends of the arcs
  \foreach \a/\b/\c/\d in {0.33/0.94/0.474/1.353, 0.568/0.16/0.818/0.23, 0.47/-0.658/0.676/-0.948, 0.148/-1.109/0.214/-1.596}{
    \draw[slate, line width=0.5pt] (\a,\b)--(\c,\d);
    \draw[slate, line width=0.5pt] (6.4-\a,\b)--(6.4-\c,\d);}
  % two paths between the same two arcs
  \draw[hot] (0.536,0.9)--(2.2,1.0)--(4.2,1.0)--(5.864,0.9);
  \draw[hot] (0.654,0.5)--(2.2,0.5)--(4.2,0.5)--(5.746,0.5);
  \draw[hot] (0.654,0.5) arc [start angle=20.9, end angle=40, x radius=0.7, y radius=1.4];
  \draw[hot] (5.746,0.5) arc [start angle=159.1, end angle=140, x radius=0.7, y radius=1.4];
  % other links
  \draw[edge] (0.684,-0.3)--(5.716,-0.3);
  \draw[edge] (0.49,-1.0)--(2.4,-1.05)--(4.0,-1.05)--(5.91,-1.0) (3.2,-1.05)--(3.2,-1.55);
  \draw[edge] (2.4,-1.05)--(3.2,-1.05)--(4.0,-1.05);
  \foreach \p in {(0.536,0.9),(2.2,1.0),(4.2,1.0),(5.864,0.9),(0.654,0.5),(2.2,0.5),(4.2,0.5),(5.746,0.5),
                  (0.684,-0.3),(5.716,-0.3),(0.49,-1.0),(2.4,-1.05),(3.2,-1.05),(4.0,-1.05),(5.91,-1.0)}
    \node[vtx] at \p {};
  \node[note, text=warm, anchor=west] at (8.1,0.75) {two paths between the same two arcs:};
  \node[note, text=warm, anchor=west] at (8.1,0.35) {a cycle with fewer than $D$ edges};
  \node[note, anchor=west] at (8.1,-0.3) {a link that is an edge};
  \node[note, anchor=west] at (8.1,-1.05) {a link that is a tree};
\end{tikzpicture}
\caption{Step~2 in the proof of Proposition~\ref{prop:loggirth}. Each cycle is divided into arcs with at most $\lfloor D/8\rfloor$
vertices. A link that is an edge or a tree gives a path with at most $t$ edges from $C$ to
$C'$. Two of these paths between the same arc of $C$ and the same arc of $C'$ would close a
cycle of $J$ with fewer than $D$ edges.}
\label{fig:twocycles}
\end{figure}

\begin{proof}
All estimates below hold for large $r$, uniformly in the other data. The family contains the
loops and the pairs of parallel edges of $\Gamma-P$, because these are cycles of length at most
two. By its maximality, $J=\Gamma-U$ is a simple graph of girth greater than $D$. We have
$s\ge(D-2)/2\ge x/32-1$, so $L^{2}2^{-s}\le2x^{2}y\,r^{-1/32}\le\eta$. With~\eqref{eq:load} this
is~(i). Recall that $|P|\le2r/x^{2}$ and that the number $c$ of components of $\Gamma[P]$ is at
most $x/100$.

\step{Step 1: components with a cycle.} Let $A$ and $B$ be disjoint regions of $\Gamma$, and let
$t\ge0$ be an integer. We claim that at most $5\cdot2^{t}$ components $Y$ of $\Gamma-(A\cup B)$
have all of the following properties: $Y$ has an edge to $A$ and an edge to $B$; $Y$ avoids $P$;
$Y$ contains a cycle; $|\delta Y|\le t$. Let $Y_0,Y_1,\dots,Y_n$ be such components. For
$1\le i<j\le n$ the graph $\Gamma-Y_i-Y_j$ is connected. Indeed, it contains the connected set
$A\cup B\cup Y_0$, and every other component of $\Gamma-(A\cup B)$ is connected and has an edge to
$A\cup B$. No edge joins $Y_i$ and $Y_j$. So the regions $Y_i$ and $Y_j$ have exactly one link.
Let $E_i$ be the event that $\cO_{Y_i}$ holds and $\xi\cap E(\Gamma[Y_i])$ is not empty. By
Lemma~\ref{lem:link}, the events $E_i$ and $E_j$ are independent, and
$\PP(E_i)\ge2^{-t}\cdot\half$. If $n=0$, the bound holds. Otherwise let $N$ be the number of the
events $E_1,\dots,E_n$ that occur, and put $\nu=\EE N\ge n\,2^{-t-1}>0$. Then
$\EE N^{2}\le\nu^{2}+\nu$, and the Cauchy--Schwarz inequality gives
$\PP(N=0)\le1-\nu^{2}/\EE N^{2}\le1/(1+\nu)$. If $\xi\cap E(\Gamma-P)$ is a forest, then no $E_i$
occurs, because $Y_i$ avoids $P$ and a nonempty even subgraph is not a forest. So
$\frac25\le q(\Gamma,P)\le1/(1+\nu)$. Hence $\nu\le\frac32$ and $n\le3\cdot2^{t}$, and there are
at most $3\cdot2^{t}+1\le5\cdot2^{t}$ such components.

\step{Step 2: the links of two cycles.} Put $t=\lfloor y/2\rfloor$; then $1\le t\le D/4$. Let
$C,C'\in\mathcal C$ be disjoint, and put $U'=V(C)\cup V(C')$. Every vertex of $U'$ lies on exactly
one edge outside its own cycle, so at most $2L$ edges leave $U'$. At most $c$ components of
$\Gamma-U'$ meet $P$, because $P\cap U'=\emptyset$ and so each component of $\Gamma[P]$ lies in one
of them. At most $2L/(t+1)$ components have more than $t$ edges to $U'$. Of the other components
that are joined to both cycles, at most $5\cdot2^{t}$ contain a cycle, by Step~1 with $A=V(C)$
and $B=V(C')$. The remaining components that are joined to both cycles induce trees, avoid $P$
and have at most $t$ edges to $U'$. They also avoid $S$: each cycle of the family is disjoint
from $U'$, so it lies in one component of $\Gamma-U'$, and this component contains a cycle.

Let $Y$ be one of the remaining components, with $d_Y\le t$ edges to $U'$. All its vertices have
degree $3$, so $3|Y|=2(|Y|-1)+d_Y$, that is, $|Y|=d_Y-2$. So there is a path from $C$ to $C'$
through $Y$ with at most $t$ edges (Figure~\ref{fig:twocycles}). For a link that is an edge,
take this edge as the path. These
paths lie in $J$. They are pairwise disjoint, because every vertex of $U'$ lies on exactly one
edge outside its own cycle. Divide each of $C$ and $C'$ into arcs of at most $\lfloor D/8\rfloor$
consecutive vertices. At most $17L/D$ arcs are necessary for each cycle, because
$\lfloor D/8\rfloor\ge D/16$ and $L\ge D$. If two of the paths had their ends in the same arc of
$C$ and in the same arc of $C'$, they would form, together with parts of these two arcs, a cycle
of $J$ of length at most $2t+2\lfloor D/8\rfloor<D$. So there are at most $(17L/D)^{2}$ such
paths. Together,
\[
  \ell(C,C')\;\le\;b:=c+\frac{2L}{t+1}+5\cdot2^{t}+289\Bigl(\frac LD\Bigr)^{2}
  \;\le\;\frac{x}{100}+\frac{4x}{\sqrt y}+5\sqrt x+289\cdot256\,y\;\le\;\frac{x}{90}.
\]

\step{Step 3: the second moment.} Call a disjoint pair $C,C'\in\mathcal{C}$ \emph{exceptional} if
$\ell(C,C')\ge3$. Fix $C$, and let $\mathcal{E}$ be a maximal family of pairwise disjoint
exceptional partners of $C$. By Lemma~\ref{lem:exc}, $|\mathcal{E}|\le|\delta V(C)|\le L$. So the
union of their vertex sets has at most $L^{2}$ vertices, and every exceptional partner of $C$
meets this union. By~\eqref{eq:load}, the sum of $a_{C'}$ over the exceptional partners $C'$ of $C$
is at most $L^{3}2^{-s}$. With~\eqref{eq:pair} and Step~2,
\[
  \sum_{C'\ \mathrm{exceptional}}\EE(X_CX_{C'})\;\le\;a_C\,2^{b}L^{3}\,2^{-s}
  \;\le\;a_C\,2^{x/90-x/32+1}\bigl(x\sqrt y\bigr)^{3}\;\le\;a_C\,\eta.
\]
By~\eqref{eq:pair}, a disjoint pair with at most one link has $\EE(X_CX_{C'})\le a_Ca_{C'}$, and
a pair with two links has $\EE(X_CX_{C'})=2a_Ca_{C'}$. Intersecting pairs contribute nothing, and
$\EE X_C^{2}\le\EE X_C=a_C$. The sum over all ordered pairs is~(iii).

\step{Step 4: the weight.} A count of edge ends gives
$3|V(\Gamma)|-3|P|\le2|E(\Gamma)|\le3|V(\Gamma)|$, and $|E(\Gamma)|=|V(\Gamma)|+r-1$. So
$2r-2\le|V(\Gamma)|\le2r-2+3|P|$. Hence the number $v$ of vertices of $J$ satisfies
$r\le v\le3r$, and $d_J\le3(|P|+|S|)\le9r/x^{2}$. By Lemma~\ref{lem:weight} with
$\ell_0=\lceil6x\rceil$,
\[
  \bigl(1+\tfrac32\,\eta\bigr)\,\mu\;\ge\;
  \Bigl(1-\frac{5\sqrt y}{x}-\frac{9}{r^{2}}\Bigr)\sum_{\ell}\frac1{2\ell},
\]
where the sum is over the even $\ell$ with $6x\le\ell\le L$. This sum is at least
$\frac14\ln\bigl(L/(6x+2)\bigr)\ge\frac14\ln\bigl(L/(6x)\bigr)-1/(12x)$, and it is at most $\ln y$.
This gives~(ii).
\end{proof}

\section{Proof of Theorem~\ref{thm:forest}}\label{sec:proof}

Let $\Gamma$ and $P$ be as in the theorem, let $\xi$ be uniform on $\Z(\Gamma)$, and let $r$ be
large. Put
\[
  x=\log r,\qquad y=\log x,\qquad z=\log y,\qquad
  D=\lfloor x/10\rfloor,\qquad L=\lfloor x\sqrt y\rfloor,\qquad\varepsilon=20/z .
\]
Let $\mathcal{F}$ be the event that $\xi\cap E(\Gamma-P)$ is a forest. We assume that
$\PP(\mathcal{F})>\frac12+\varepsilon$ and derive a contradiction.

\step{Few short cycles.} Let $\mathcal{Q}$ be a maximal family of pairwise disjoint cycles of
$\Gamma-P$ of length at most $D$, and let $S$ be the union of their vertex sets. For
$Q\in\mathcal{Q}$ the set $V(Q)$ is a cyclic region of $\Gamma$ with
$|\delta V(Q)|\le|Q|\le D$, because its vertices have degree $3$. On $\mathcal{F}$ the set
$\xi\cap E(\Gamma[V(Q)])$ is a forest for every $Q$. So Lemma~\ref{lem:smallcuts} gives
$|\mathcal{Q}|<(D+1)2^{D-1}/\varepsilon\le\sqrt r$. Hence $|S|\le D\sqrt r\le r/x^{2}$.

\step{The cycles of moderate length.} Now Proposition~\ref{prop:loggirth} applies, with the
variable $Z$ of Section~\ref{sec:events} for $U=P\cup S$. As $L/(6x)\ge\sqrt y/7$ and
$W_2\le\mu^{2}$, it gives
\[
  \mu\;\ge\;\frac14\ln\frac{\sqrt y}{7}-\frac1{\sqrt x}\;\ge\;\frac{\ln y}{10},
  \qquad \EE Z^{2}\;\le\;2\mu^{2}+2\mu .
\]
On $\mathcal{F}$ no cycle of $J\subseteq\Gamma-P$ is contained in $\xi$, so $Z=0$. By
Lemma~\ref{lem:second},
\[
  \PP(\mathcal{F})\;\le\;\PP(Z=0)\;\le\;\frac12+\frac1{2\mu}\;\le\;\frac12+\frac{5}{\ln y}
  \;<\;\frac12+\frac{20}{z},
\]
since $\ln y=z\ln2$. This is a contradiction, and Theorem~\ref{thm:forest} is proved. \qed

\section{Simple graphs}\label{sec:simple}

We prove Corollary~\ref{cor:simple}. Let $G$ and $W$ be as in the corollary, with $m=|E(W)|$
and $c=\comp(W)$. We construct a multigraph $\Gamma$ and a set $P$ as in
Theorem~\ref{thm:forest}, with $p(G,W)\le q(\Gamma,P)$. The construction has two steps
(Figure~\ref{fig:split}).

\begin{figure}[ht]
\centering
\begin{tikzpicture}[x=1cm,y=1cm]
  % before
  \coordinate (v) at (0,0);
  \draw[cool] (v)--(-1.05,0.7) (v)--(-1.05,-0.7);
  \draw[edge] (v)--(0.15,1.05) (v)--(1.05,0.6) (v)--(1.05,-0.6);
  \node[vtx] at (v) {};
  \node[note, text=accent] at (-1.45,0) {$W$};
  \node[note] at (0.05,-0.4) {$v$};
  \draw[slate, line width=0.6pt, -{Stealth[length=4.5pt]}] (2.1,0)--(3.3,0);
  % after
  \coordinate (v1) at (5.4,0);
  \coordinate (v2) at (6.7,0);
  \coordinate (v3) at (8.0,0);
  \draw[draw=accent, fill=accent!7, line width=0.6pt] (v1) circle [radius=0.3];
  \draw[cool] (v1)--(4.35,0.7) (v1)--(4.35,-0.7);
  \draw[hot] (v1)--(v2)--(v3);
  \draw[edge] (v2)--(6.7,1.05) (v3)--(9.05,0.6) (v3)--(9.05,-0.6);
  \foreach \p in {v1,v2,v3} \node[vtx] at (\p) {};
  \node[note, text=accent] at (3.95,0) {$W$};
  \node[note] at (5.4,-0.55) {$v^{1}$};
  \node[note] at (6.7,-0.4) {$v^{2}$};
  \node[note] at (8.0,-0.4) {$v^{3}$};
  \node[note, text=accent, anchor=west] at (9.6,0.45) {$v^{1}$ is put into $P$};
  \node[note, text=warm, anchor=west] at (9.6,-0.05) {new path edges};
\end{tikzpicture}
\caption{Splitting a vertex $v$ of degree five that lies on two edges of $W$. The edges of $W$
come first, so they end at $v^{1}$.}
\label{fig:split}
\end{figure}

\step{Splitting.} Replace every vertex $v$ of degree $d\ge4$ by a path $v^{1}v^{2}\cdots v^{d-2}$.
Distribute the $d$ edges at $v$ along this path in some order: two to $v^{1}$, one to each of
$v^{2},\dots,v^{d-3}$, and two to $v^{d-2}$. If $v\in V(W)$, let the edges of $W$ at $v$ come first
in this order, so that they end in an initial segment $v^{1},\dots,v^{i}$ of the path with
$i<\deg_W(v)$. Let $G'$ be the new graph; its degrees are at most $3$. Let $P$ consist of these
initial segments and of the vertices of $W$ that were not split. Then $|P|\le\sum_v\deg_W(v)=2m$,
every edge of $W$ joins two vertices of $P$, and $G'[P]$ has at most $c$ components.

Contracting the new path edges maps $\Z(G')$ bijectively onto $\Z(G)$. So a uniform even subgraph
$\xi'$ of $G'$ gives the uniform even subgraph $\xi=\xi'\cap E(G)$ of $G$. Every edge of $G'-P$
is a new path edge or an edge of $E(G)\sm E(W)$. If $\xi\sm E(W)$ is a forest, then so is its
union with the set of all new path edges, because contracting these edges, which form a forest,
gives back $\xi\sm E(W)$. Hence $p(G,W)\le q(G',P)$.

\step{Cleaning.} As long as possible, delete a vertex of degree one outside $P$, or replace a vertex
of degree two outside $P$ and its two edges by a single edge. An even subgraph avoids a pendant edge
and contains both or none of the two edges at a vertex of degree two. So these operations induce
bijections of the spaces of even subgraphs. Let $\Gamma$ be the final multigraph. Each edge of
$\Gamma-P$ is a path of $G'$ all of whose vertices lie outside $P$. So if an even subgraph of
$\Gamma$ contains a cycle of $\Gamma-P$, then the corresponding even subgraph of $G'$ contains a
cycle of $G'-P$. Hence $q(G',P)\le q(\Gamma,P)$.

\step{The estimate.} The multigraph $\Gamma$ is connected and subcubic, and it has rank $r$.
Edges with both ends in $P$ are never removed, so the number of components of the graph induced
on $P$ does not increase. The vertices outside $P$ now have degree exactly $3$. As
$|P|\le2m\le2r/\log^{2}r$ and $c\le(\log r)/100$, Theorem~\ref{thm:forest} gives
$p(G,W)\le q(\Gamma,P)\le\half+20/\lgk{3}r$. \qed

\begin{thebibliography}{9}\small

\bibitem{Bass} H.~Bass, The Ihara--Selberg zeta function of a tree lattice, \emph{Internat.\
J.~Math.}\ \textbf{3} (1992), 717--797.

\bibitem{Bloom} T.~F.~Bloom, Erd\H{o}s Problem \#1016, \url{https://www.erdosproblems.com/1016},
accessed November 2026.

\bibitem{ES} R.~C.~Entringer and P.~J.~Slater, On the maximum number of cycles in a graph,
\emph{Ars Combin.}\ \textbf{11} (1981), 289--294.

\bibitem{HK} P.~Horak and D.~Kim, Connected cubic graphs with the maximum number of perfect
matchings, \emph{J.~Graph Theory} \textbf{99} (2022), 671--690. arXiv:2006.13459.

\bibitem{KNT} KNT, The minimum number of edges in a pancyclic graph, manuscript, September 2026,
\url{https://jehlp.net/erdos1016/}.

\bibitem{KS} M.~Kotani and T.~Sunada, Zeta functions of finite graphs, \emph{J.~Math.\ Sci.\ Univ.\
Tokyo} \textbf{7} (2000), 7--25.

\bibitem{Max} KNT, Astra and Fable~5.2, The maximum number of cycles in a graph of large rank,
companion manuscript, November 2026.

\bibitem{Pan} KNT, Astra and Fable~5.2, The minimum number of edges in a pancyclic graph: a
short proof, companion manuscript, November 2026.

\bibitem{Repo} KNT, Astra and Fable~5.2, Lean~4 formalisation of the forest estimate and of the
results of~\cite{Pan} and~\cite{Max}, \url{https://github.com/JWKNT/erdos1016}.

\end{thebibliography}

\end{document}
